\section*{Bab \ref{ch:b3:gfd-foundations} --- Dinamika Fluida Geofisika}

\begin{solution}{exo:b3:gfd-foundations:1}
(a)~$f = 2\Omega\sin\varphi$: \num{7.292e-5}, \num{1.031e-4}, dan
\qty{1.263e-4}{s^{-1}}; $2\pi/f$: 23.93, 16.92, dan \qty{13.82}{h}.
(b)~$U/f$: 2.74, 1.94, dan \qty{1.58}{km}. (c)~Pada \ang{30}, $\sin\varphi =
1/2$ dan $2\pi/f = 2\pi/\Omega$: periode inersia adalah hari sideral,
\qty{23.93}{h}.
\end{solution}

\begin{solution}{exo:b3:gfd-foundations:2}
Dengan $f = \qty{1.031e-4}{s^{-1}}$: pusaran debu \num{9.7e3}, angin laut 2.4,
depresi 0.15, gyre \num{2.4e-4}. Depresi dan gyre mendekati
keseimbangan geostrofik; pusaran debu berputar ke kedua arah, tidak menyadari
rotasi Bumi, dan angin laut merasakannya hanya selama sehari.
\end{solution}

\begin{solution}{exo:b3:gfd-foundations:3}
$\partial u_g/\partial z = (g/fT)\,|\partial T/\partial y| = 9.81\times
\num{5e-6}/(\num{1.031e-4}\times250) = \qty{1.90e-3}{s^{-1}}$; sepanjang
\qty{10}{km} angin barat meningkat sebesar \qty{19}{m/s}.
\end{solution}

\begin{solution}{exo:b3:gfd-foundations:4}
$E = \num{50e-8}$: $N = \qty{2.21e-3}{s^{-1}}$, periode \qty{47}{min},
\num{1.27} siklus per jam. $E = \num{1000e-8}$: $N = \qty{9.90e-3}{s^{-1}}$,
\qty{10.6}{min}, 5.7 siklus per jam. $E = \num{e-8}$: $N =
\qty{3.13e-4}{s^{-1}}$, \qty{5.6}{h}, 0.18 siklus per jam.
\end{solution}

\begin{solution}{exo:b3:gfd-foundations:5}
(a)~$W \sim UH/L = \qty{0.1}{m/s}$ dan $\mathrm Dw/\mathrm Dt \sim UW/L =
\qty{e-6}{m/s^2}$, \num{1.0e-7} kali $g$. Suku tekanan, $fUL/H =
\qty{0.103}{m/s^2}$, lebih besar \num{1e5} kali: rasionya, \num{9.7e-6}, adalah
$\mathrm{Ro}(H/L)^2$ dengan $\mathrm{Ro} = 0.097$. (b)~$\tilde fu =
2\Omega\cos\ang{45}\times10 = \qty{1.03e-3}{m/s^2}$, \num{1.05e-4} kali $g$.
\end{solution}

\begin{solution}{exo:b3:gfd-foundations:6}
$fr = \qty{30}{m/s}$. (a)~$\mathrm{Ro}_g = 0.667$, $V = 40/(1 + \sqrt{3.667})
= \qty{13.7}{m/s}$, sepertiga lebih lambat daripada geostrofi. (b)~$\mathrm{Ro}_g =
0.167$, $V = 10/(1 + \sqrt{0.333}) = \qty{6.34}{m/s}$. (c)~$V_g \le fr/4 =
\qty{7.5}{m/s}$, dan kemudian $V = fr/2 = \qty{15}{m/s}$, dua kali
kecepatan geostrofik.
\end{solution}

\begin{solution}{exo:b3:gfd-foundations:7}
(a)~$\Omega = 2\pi/6 = \qty{1.047}{rad/s}$, $f = \qty{2.094}{s^{-1}}$.
$\mathrm{Ro} = 0.01/(2.094\times0.05) = 0.095$ dan $\mathrm{Ek} =
\num{e-6}/(2.094\times0.04) = \num{1.19e-5}$: keduanya kecil, sehingga aliran mematuhi
\omterm{thm:b3:gfd-foundations:taylor-proudman}{teorema Taylor--Proudman} dan \omterm{thm:b3:gfd-foundations:taylor-proudman}{kolom Taylor} terbentuk. (b)~$0.2/\sqrt{\num{e-6}
\times1.047} = \qty{195}{s}$, \qty{3.3}{min}, 33 putaran meja.
\end{solution}

\begin{solution}{exo:b3:gfd-foundations:8}
$\partial\bar T/\partial y = -f\,\Delta u/(R\ln(850/250))$, dengan
$\ln(850/250) = 1.224$. Pada \ang{30} LU, $-\num{7.29e-5}\times36.12/(286.9
\times1.224) = \qty{-7.50}{K}$ per \qty{1000}{km}; pada \ang{45} LU,
$-\num{1.031e-4}\times17.18/351.1 = \qty{-5.05}{K}$ per \qty{1000}{km}. Geser
sebanding dengan gradien dibagi $f$: gradiennya 1.5
kali lebih besar pada \ang{30} dan $f$ 1.4 kali lebih kecil, sehingga geser 2.1 kali
lebih besar.
\end{solution}

\begin{solution}{exo:b3:gfd-foundations:9}
(a)~$L_R = 0.0113\times11\,000/\num{1.031e-4} = \qty{1206}{km}$;
$0.00357\times1000/\num{9.47e-5} = \qty{37.7}{km}$, 32 kali lebih kecil.
(b)~$\beta = 2\Omega\cos\varphi/a$: \num{1.62e-11} dan
\qty{1.74e-11}{m^{-1}.s^{-1}}. $\beta L_R/f = 0.19$ untuk atmosfer, 0.007
untuk samudra.
\end{solution}

\begin{solution}{exo:b3:gfd-foundations:10}
(a)~$\mathrm dw/\mathrm dt = -(r + \mathrm if)w$, sehingga $w = U\mathrm
e^{-rt}\mathrm e^{-\mathrm ift}$: kecepatan berputar dengan laju $f$ sementara
magnitudonya meluruh sebagai $\mathrm e^{-rt}$. Posisi $x + \mathrm iy =
U(1 - \mathrm e^{-(r + \mathrm if)t})/(r + \mathrm if)$ berspiral ke
titik $U/(r + \mathrm if)$. (b)~$x + \mathrm iy = U(1 - \mathrm e^{-(r +
\mathrm if)t})/(r + \mathrm if) + u_0t$: spiral itu tertarik memanjang menjadi gelung.
(c)~$0.2\,\mathrm e^{-1} = \qty{0.074}{m/s}$; pusat telah bergeser $u_0t =
0.05\times4\times86\,400 = \qty{17.3}{km}$ ke timur.
\end{solution}

\begin{solution}{exo:b3:gfd-foundations:11}
$f = \qty{9.37e-5}{s^{-1}}$; $\partial u/\partial z = (g/f\rho_0)\,\partial
\rho/\partial y = 9.81\times10^{-5}/(\num{9.37e-5}\times1025) =
\qty{1.02e-3}{s^{-1}}$. Dengan $u = 0$ pada \qty{1000}{m} arus permukaan adalah
\qty{1.02}{m/s}, ke timur: di belahan bumi utara arus mengalir dengan
air yang lebih ringan di sebelah kanannya.
\end{solution}

\begin{solution}{exo:b3:gfd-foundations:12}
(a)~$N^2 = (g/\rho)\,\Delta\sigma_\theta/\Delta z$. 0--\qty{100}{m}:
$9.81\times0.668/(1026.4\times100)$, $N = \qty{7.99e-3}{s^{-1}}$, 4.6
siklus per jam; 100--\qty{500}{m}: \qty{3.02e-3}{s^{-1}}, 1.7 siklus per
jam; 500--\qty{1000}{m}: \qty{3.24e-3}{s^{-1}}, 1.9 siklus per jam.
(b)~\qty{100}{m} teratas: lapisan ini mengandung air permukaan yang hangat di atas
piknoklin musiman. Antara 100 dan \qty{500}{m} air lebih
seragam daripada di bawahnya. (c)~$E = N^2/g$ adalah \num{650e-8}, \num{93e-8}, dan
\qty{107e-8}{m^{-1}}, semuanya berada dalam kisaran samudra bagian atas.
\end{solution}

\begin{solution}{pb:b3:gfd-foundations:1}
\textbf{1.} $f = 2\times\num{7.292115e-5}\sin\ang{48} =
\qty{1.084e-4}{s^{-1}}$; $2\pi/f = \qty{16.1}{h}$.

\textbf{2.} Bak mandi \num{3.1e3}; tangki ($f = 4\pi/10 = \qty{1.257}{s^{-1}}$)
0.020; angin laut 2.3; depresi 0.14; gyre \num{2.3e-4}.

\textbf{3.} Bak mandi sama sekali tidak merasakan rotasi Bumi; angin laut
merasakannya hanya setelah beberapa jam. Tangki, depresi, dan gyre memiliki
\omterm{def:b3:gfd-foundations:numbers}{bilangan Rossby} kecil dan mendekati geostrofik.

\textbf{4.} Tangki $\num{e-6}/(1.257\times0.04) = \num{2.0e-5}$; depresi
$10/(\num{1.084e-4}\times10^8) = \num{9.2e-4}$; gyre
$\num{e-3}/(\num{1.084e-4}\times10^6) = \num{9.2e-6}$.

\begingroup\emergencystretch=3em\relax \textbf{5.} $\delta/H = \sqrt{2\nu/(fH^2)} = \sqrt{2\,\mathrm{Ek}}$. Tangki
$\sqrt{2\times\num{e-6}/1.257} = \qty{1.3}{mm}$; atmosfer $\sqrt{20/\num{1.084e-4}}
= \qty{430}{m}$.
\par\endgroup
\textbf{6.} Gesekan bekerja pada lapisan yang tebalnya beberapa perseribu kedalaman di
tangki dan samudra, beberapa perseratus di atmosfer; di atasnya aliran
hampir tanpa gesekan.

\textbf{7.} $\rho = p/(RT) = 101\,290/(286.9\times288.1) =
\qty{1.225}{kg/m^3}$.

\textbf{8.} $u_g = (1/\rho f)\,|\partial p/\partial y| = (400/300\,000)/(1.225
\times\num{1.084e-4}) = \qty{10.0}{m/s}$, dari barat, dengan tekanan rendah
berada di utara, di sebelah kiri.

\textbf{9.} $\mathrm{Ro}_g = 10.04/(\num{1.084e-4}\times600\,000) = 0.154$;
$V = 2\times10.04/(1 + \sqrt{1.618}) = \qty{8.84}{m/s}$, \qty{12}{\percent}
lebih lambat.

\textbf{10.} $V = 2\times10.04/(1 + \sqrt{0.382}) = \qty{12.4}{m/s}$. Nilai
$V_g$ terbesar adalah $fr/4 = \qty{16.3}{m/s}$, gradien $\rho fV_g =
\qty{2.16e-3}{Pa/m}$, \qty{6.5}{hPa} per \qty{300}{km}.

\textbf{11.} $g/c_p = \qty{9.77}{K/km}$. $N^2 = (9.81/255.7)\times(9.77 -
6.49)\times10^{-3} = \qty{1.26e-4}{s^{-2}}$, $N = \qty{0.0112}{s^{-1}}$, dengan
periode \qty{9.3}{min}.

\textbf{12.} Dengan $\Gamma = 0$, $N^2 = (g/T)(g/c_p) = g^2/(c_pT) =
\qty{4.43e-4}{s^{-2}}$, $N = \qty{0.0210}{s^{-1}}$, 1.9 kali
nilai troposfer.

\textbf{13.} $L_R = 0.0112\times10^4/\num{1.084e-4} = \qty{1035}{km}$, kira-kira
sama dengan ukuran depresi.

\textbf{14.} $g' = \qty{0.0147}{m/s^2}$; $\sqrt{g'h} = \qty{0.54}{m/s}$;
$L_R = 0.542/\num{1.084e-4} = \qty{5.0}{km}$.

\textbf{15.} Ya: danau selebar \qty{12}{km}, lebih dari dua kali
\omterm{def:b3:gfd-foundations:rossby-radius}{jari-jari deformasi Rossby-nya}, sehingga gelombang dan arus internal di dalamnya dibelokkan oleh
gaya Coriolis.

\textbf{16.} $(249.4 - 253.0)/600 = \qty{-6.0}{K}$ per \qty{1000}{km}.

\textbf{17.} $\ln(1013/265) = 1.3409$. Selatan:
$286.9\times253.0\times1.3409/9.81 = \qty{9922}{m}$; utara \qty{9781}{m}.
Permukaan \qty{265}{hPa} lebih rendah \qty{141}{m} di utara, dengan kemiringan
\num{2.35e-4} relatif terhadap permukaan bawah.

\textbf{18.} $\Delta u = (g/f)\times\text{kemiringan} = 9.81\times\num{2.35e-4}/
\num{1.084e-4} = \qty{21.3}{m/s}$.

\textbf{19.} $(R/f)\ln(p_b/p_t)\,|\partial\bar T/\partial y| =
(286.9/\num{1.084e-4})\times1.3409\times\num{6.0e-6} = \qty{21.3}{m/s}$, hasilnya
sama: kedua perhitungan itu sebenarnya satu.

\textbf{20.} \omterm{def:b3:gfd-foundations:baroclinic}{Baroklinik}: suhu, dan karena itu massa jenis, berubah sepanjang
permukaan tekanan. \omterm{def:b3:gfd-foundations:baroclinic}{Barotropik}, angin tidak akan berubah terhadap ketinggian:
\qty{10.0}{m/s}.

\textbf{21.} $(g/f\bar T)\,|\partial T/\partial y|\,\Delta z = 9.81\times
\num{6.0e-6}\times10^4/(\num{1.084e-4}\times255.7) = \qty{21.2}{m/s}$,
dalam \qty{0.3}{\percent}.

\textbf{22.} $10.04 + (286.9/\num{1.084e-4})\times\num{6.0e-6}\times
\ln(1013/541) = 10.04 + 9.96 = \qty{20.0}{m/s}$.

\textbf{23.} $\partial u/\partial z = 21.3/10^4 = \qty{2.13e-3}{s^{-1}}$;
$\mathrm{Ri} = \num{1.26e-4}/\num{4.54e-6} = 28$, jauh di atas 0.25: geser
tidak membuat troposfer bebas menjadi turbulen.

\textbf{24.} Angin Orrin, \qty{31}{m/s}, lebih kuat daripada rata-rata zonal
pada \ang{45} LU dan lebih lemah daripada pada \ang{30} LU. Arus jet front kutub
bergeser ke utara dan selatan; pada sembarang saat gradien dan anginnya
terkonsentrasi di sebuah pita, dan rata-rata sepanjang lingkaran lintang menyebarkannya.

\textbf{25.} $10.0 + 21.3$: angin geostrofik pada \qty{10}{km} di atas
stepa adalah \textbf{\qty[text-series-to-math=true]{31}{m/s}}, dari barat
(\qty{31.3}{m/s} menurut perhitungan).
\end{solution}
